% ============================================================
% Part III — Desktop Section — M1 (The Two Literatures) — DRAFT v1 (rev 2, 2026-07-20)
% Mode: Stage 3 guided drafting, parallel session (module-by-module cadence)
% Blueprint: PART-III-DESKTOP-SECTION-OUTLINE-v1-2026-07-15.md, M1 as revised
%   2026-07-19 (M1.2+M1.3 merged; M1.5 deferred to the boredom-experiment
%   section; four movements total). Sectioning level provisional pending
%   Open decision #1 (arc placement).
% Conventions: MLA parenthetical cites (abbreviated titles where an author
%   set has multiple works; bare page for same-work repeat within a
%   paragraph); footnotes substantive only; first person for King & Salvo
%   (full MLA cite + convention footnote lives in the M0 draft);
%   NO abstract quotations; no sentence-initial "And"; "record" banned as a
%   name for the evidence (use "the survey data"); no metatextual signposting.
% REV 2 (2026-07-20, per user notes + cross-session verification): fidelity
%   DE-EMPHASIZED throughout M1.1 — the psychological-determination verdict
%   (DA-03), M0.3's locked claim (world-disclosedness does not require a
%   headset), and Dubovi's own Table 3 (only Involved/Control correlates with
%   learning) all demote representational fidelity to one tool among several;
%   interaction/task now leads wherever the pair is named. Dubovi quote kept
%   as expansion-and-disagreement per user; vague "half of that sentence"
%   framing replaced with plain statement of the two clauses.
% Quotes: PDF-verified per sources/vle-pedagogy-lit-quote-bank.md (ROUND 2
%   dispositions folded in). Pages empirical from headers/footers.
%   NB Dalgarno & Lee 2012: no printed footers in PDF; pages inferred from
%   the proceedings range (236-245) — bank §0.
% Batch 1 (2026-07-19, rev 2 2026-07-20): M1.1. Batch 2 (2026-07-20): merged
%   M1.2/3 below.
%
% WORKS CITED entries introduced by this module (for later assembly):
%   Dalgarno, Barney, and Mark J. W. Lee. "What Are the Learning Affordances
%     of 3-D Virtual Environments?" British Journal of Educational
%     Technology, vol. 41, no. 1, 2010, pp. 10-32.
%   ---. "Exploring the Relationship between Afforded Learning Tasks and
%     Learning Benefits in 3D Virtual Learning Environments." Future
%     Challenges, Sustainable Futures: Proceedings ascilite Wellington 2012,
%     ascilite, 2012, pp. 236-45.
%   Dubovi, Ilana, et al. "Now I Know How! The Learning Process of Medication
%     Administration among Nursing Students with Non-Immersive Desktop
%     Virtual Reality Simulation." Computers & Education, vol. 113, 2017,
%     pp. 16-27.
%   Fowler, Chris. "Virtual Reality and Learning: Where Is the Pedagogy?"
%     British Journal of Educational Technology, vol. 46, no. 2, 2015,
%     pp. 412-22.
%   Makransky, Guido, and Gustav Bøg Petersen. "Investigating the Process of
%     Learning with Desktop Virtual Reality: A Structural Equation Modeling
%     Approach." Computers & Education, vol. 134, 2019, pp. 15-30.
%   Merchant, Zahira, et al. "The Learner Characteristics, Features of
%     Desktop 3D Virtual Reality Environments, and College Chemistry
%     Instruction: A Structural Equation Modeling Analysis." Computers &
%     Education, vol. 59, 2012, pp. 551-68.
%   Mikropoulos, T. A., and A. Natsis. "Educational Virtual Environments: A
%     Ten-Year Review of Empirical Research (1999-2009)." Computers &
%     Education, vol. 56, 2011, pp. 769-80.  [initials kept until verified]
% ============================================================

\section{The Two Literatures}

% --- M1.1 Desktop 3D as a modality in its own right (rev 2, 2026-07-20)

The desktop route has a research literature of its own, and that literature's
founding model already refuses the reading a headset-first field invites: the
desktop build as a lesser headset. Dalgarno and Lee's model of learning in
three-dimensional virtual environments locates what is pedagogically
distinctive about the medium in what an environment lets its learner do and
see---learner interaction first, representational fidelity alongside it---and
treats presence and co-presence not as endowments of hardware but as
consequences of those characteristics: in the model's words, ``a consequence
of the representational fidelity and learner interactions facilitated by the
environment'' (``Learning Affordances'' 14). The affordances the model
derives are, accordingly, task-possibilities in Gibson's sense---things an
environment makes available to an agent to do---not display achievements.
The model arrived carrying its authors' own verdict on the field's evidence:
``[m]uch of what has been published about the educational uses of 3-D
technologies is largely `show-and-tell', presenting only anecdotal evidence
or personal impressions that cannot be usefully generalised beyond the local
context'' (23). When the same authors put the model to its own test two
years later---surveying fifty-three higher-education staff on which afforded
learning tasks had yielded which of the five predicted benefits---respondents
endorsed every benefit at nearly uniform strength regardless of task, none of
the hypothesized task-to-benefit links reached significance, and the authors
flagged the deflating explanation themselves: the uniform agreement ``could
be more a reflection of their enthusiasm and positive attitude towards the
use of 3D VLEs than an evidence-based assessment of the actual learning
benefits'' (``Afforded Tasks'' 243). The field's central model thus stands
named beside its own inconclusive self-test.

The model's sharpest critique asked not for more data but for pedagogy: a
decade of educational-VR research had run technology-first, deriving expected
benefits from technical affordances with almost no explicit account of how
learning occurs.\footnote{The ten-year review behind the charge found
explicit learning theory almost wholly absent from the literature it
surveyed: ``[a]ll the other reviewed articles do not refer explicitly to a
learning theory'' (Mikropoulos and Natsis 775, qtd.\ in Fowler 412). Fowler's
repair works outward from intended learning outcomes to the technology that
would serve them, rather than inward from the technology's affordances
(412--13, 416--20).} What Fowler's repair adds, and what this section will
hold onto, is a curve: ``the assumption that high levels of representational
fidelity and learner interaction will result in deeper learning is
questionable. It may be argued that there are optimum levels of these
characteristics that maximise learning. Going beyond the optimum increases
costs and results in a limited or even negative return with respect to
learning benefits'' (Fowler 415). Neither characteristic improves learning in
proportion to its own increase: each has an optimum level, and beyond it,
added fidelity or added interaction yields limited or even negative returns
for learning. That caution matters most for representational fidelity,
because fidelity is the characteristic that headset hardware exists to
advance---and the one a headset-first field kept mistaking for the medium's
pedagogical substance. Our own pilot had already produced a concrete instance
of that curve. We expected that students' familiarity with VR and video games would
incline them toward the virtual-reality versions; the correlation ran the
other way: of the students who judged the plain 2D videos most
educationally beneficial, 66.7\% were weekly gamers, against 33.3\% of those
who favored the VR versions---and our interviews traced the judgment to
builds that ``do not meet the production standards of contemporary VR and
traditional video games'' (``Phenomenological Evaluation'' 389). We stated
the result then at its proper strength---a possible correlation, awaiting
validation (389)---but its direction is the one the curve predicts once
experience is allowed to move the optimum: the students most practiced in
polished, feature-rich worlds hold the highest bar, and a build that
undershoots it loses them first.

Within the desktop medium itself, the process studies converge on where the
leverage lies, and it is not the display. Makransky and Petersen, modeling
199 medical undergraduates through a desktop genetics simulation, found two
routes from the medium's features to knowledge gains---``the affective and
the cognitive paths'' (``Investigating the Process'' 24): presence feeding
motivation on one, perceived usability feeding cognitive benefit on the
other, the two converging on self-efficacy as the sole direct predictor of
learning.\footnote{Self-efficacy $\rightarrow$ learning, $\beta = 0.579$;
presence's effect reaches learning only through motivation and then
self-efficacy, never directly (``Investigating the Process'' 23--24).} In
% FLAG 2026-07-31 (from M0 Second Life footnote verification): index unit
% ped-sec-16 records the SL treatment as three graded activities over roughly
% eight weeks, not a full "semester of Second Life chemistry" — user to rule
% on rewording at the subsection-2 readthrough.
Merchant and colleagues' semester of Second Life chemistry (N\,=\,204), the
gate is usability: ``the 3D virtual reality features can support the
development of learners' spatial orientation ability, self-efficacy, and
presence only when the learners' [sic] perceive the experience as meaningful
and the system easy to use'' (Merchant et al.\ 562).\footnote{The mediation
is nearly total ($\beta = 0.956$, VR features $\rightarrow$ usability;
Merchant et al.\ 560). The design is single-condition desktop---no headset
arm---so the study licenses no claim about media sufficiency in either
direction.} Dubovi, Levy, and Dagan's nursing simulation, finally, returned
the literature's flattest correction to the hardware account of presence: on
ordinary monitors, their students' presence scores met or exceeded the norms
the construct's founding instrument had recorded for immersive
VR,\footnote{Involved/Control 5.2\,$\pm$\,0.7, Natural 4.7\,$\pm$\,1.03,
Interface Quality 5.0\,$\pm$\,1.1, against Witmer and Singer's immersive-VR
norms of 5.2\,$\pm$\,0.8, 4.1\,$\pm$\,1.1, and 4.8\,$\pm$\,1.1 respectively
(Dubovi et al.\ 23).} and they drew a two-part conclusion whose parts deserve
different fates: ``presence is not a unique attribute of the environment, but
it is induced by the fidelity of representation and the high degree of
interaction and user control'' (Dubovi et al.\ 25).

The first clause is right, and the survey data I will analyze later on
confirm it at scale:
presence is not what headsets have and screens lack. The second clause gives
fidelity a role the evidence does not sustain---not the wider field's
evidence, and not this study's own. The one component of their presence
measure that correlated with learning was involvement and control
($r = 0.211$ overall, $0.231$ on the procedure's core subscale); the
realism-adjacent components tracked nothing (25).\footnote{Their own design
principle, stated two sentences before the fidelity clause, names control
alone: ``ensuring that participants gain sufficient control of the
environment is an important condition for learning'' (Dubovi et al.\ 25).
The fidelity clause is carried by citation to prior literature, not by the
study's own data.} Presence flourishes in worlds whose fidelity is schematic
or frankly poor---the cube-built landscapes of \textit{Minecraft}, the aging
textures of \textit{World of Warcraft}---worlds no fidelity-led account
explains. What induces presence is not the exactness of the copy but the
constitution of the artifact as a world: an environment built as an ecology,
one that answers action---a world the student moves and is moved by in turn.
Fidelity is one tool for building such a world. It is not the ground of it.

% --- merged M1.2/3 Presence on screens, and its limit (batch 2, 2026-07-20)
% Quotes PDF-verified per bank §2 (Lombard & Ditton by direct PDF read —
%   OCR-LOW protocol; JCMC unpaginated -> section-name cites). Mayer et al.
%   = 2022 version-of-record (print issue 2023). Parong & Mayer 2018 pages
%   = APA Online First; pin at assembly. Additional Works Cited entries:
%   Barrett, Robin C. A., et al. "Comparing Virtual Reality, Desktop-Based
%     3D, and 2D Versions of a Category Learning Experiment." PLOS ONE,
%     vol. 17, no. 10, 2022, e0275119.
%   Calleja, Gordon. In-Game: From Immersion to Incorporation. MIT Press,
%     2011.
%   Chow, Meyrick. "Determinants of Presence in 3D Virtual Worlds: A
%     Structural Equation Modelling Analysis." [venue pinned at assembly],
%     2016.
%   Cook, David A., et al. "Technology-Enhanced Simulation for Health
%     Professions Education: A Systematic Review and Meta-Analysis." JAMA,
%     vol. 306, no. 9, 2011, pp. 978-88.
%   Kononowicz, Andrzej A., et al. "Virtual Patient Simulations in Health
%     Professions Education: Systematic Review and Meta-Analysis by the
%     Digital Health Education Collaboration." Journal of Medical Internet
%     Research, vol. 21, no. 7, 2019, e14676.
%   Lombard, Matthew, and Theresa Ditton. "At the Heart of It All: The
%     Concept of Presence." Journal of Computer-Mediated Communication,
%     vol. 3, no. 2, 1997.
%   Makransky, Guido, Thomas S. Terkildsen, and Richard E. Mayer. "Adding
%     Immersive Virtual Reality to a Science Lab Simulation Causes More
%     Presence but Less Learning." Learning and Instruction, vol. 60, 2019,
%     pp. 225-36.
%   Mayer, Richard E., Guido Makransky, and Jocelyn Parong. "The Promise
%     and Pitfalls of Learning in Immersive Virtual Reality." International
%     Journal of Human-Computer Interaction, vol. 39, no. 11, 2022,
%     pp. 2229-38.
%   Parong, Jocelyn, and Richard E. Mayer. "Learning Science in Immersive
%     Virtual Reality." Journal of Educational Psychology, vol. 110, no. 6,
%     2018, pp. 785-97.
%   ---. "Cognitive and Affective Processes for Learning Science in
%     Immersive Virtual Reality." Journal of Computer Assisted Learning,
%     vol. 37, 2021 [pages pinned at assembly].
%   Witmer, Bob G., and Michael J. Singer. "Measuring Presence in Virtual
%     Environments: A Presence Questionnaire." Presence: Teleoperators and
%     Virtual Environments, vol. 7, no. 3, 1998, pp. 225-40.

Presence itself entered the research literature as a property of persons,
not of headsets. Lombard and Ditton's founding synthesis defines it as
``the perceptual illusion of nonmediation''---an illusion that ``occurs when
a person fails to perceive or acknowledge the existence of a medium in
his/her communication environment and responds as he/she would if the medium
were not there''---and the definition is deliberately medium-general, built
to hold across television, film, and virtual reality alike (``Presence
Explicated''). The consequence is drawn in the same section: ``[b]ecause it
is a perceptual illusion, presence is a property of a person,'' varying with
a medium's form, its content, and the person attending to it (``Presence
Explicated''). The founding instrument agrees about the location. Witmer and
Singer, constructing the field's first validated presence questionnaire,
reject the view that immersion is an objective property of the display
hardware: ``immersion, like involvement and presence, is something the
individual experiences'' (227). Their paper is equally plain about what had
not been shown: factors that increase presence were expected to increase
learning, ``but as yet we have no direct evidence to support our
contention'' (239). Two decades of determinants research kept resolving the
question user-side---in Chow's structural model of a desktop virtual world,
computer self-efficacy, the student's confidence with the machine, was the
strongest predictor of presence, ahead of any property of the display
(11--13).\footnote{Chow's model is scoped to self-presence rather than
spatial presence, and his own design implication follows the data: invest in
orientation and user confidence rather than in richer rendering alone
(12--13).} Across research traditions that rarely cite one another, the
convergence is the same: presence is achieved by persons and occasioned by
worlds; no hardware tier owns it.

The media-comparison studies of the past two decades tested the direct
question: does the added presence of a more immersive medium produce more
learning? Across content-matched designs, the answer was no.
Makransky, Terkildsen, and Mayer ported an
identical laboratory simulation from desktop monitor to head-mounted
display: presence rose sharply in the headset (d\,=\,1.30 in the first
lesson, 2.20 in the second), tested learning fell, and EEG measures located
the mechanism in cognitive overload (``Adding Immersive'' 232--33).
\footnote{The dissociation also holds inside a single desktop condition: in
Merchant and colleagues' chemistry model, usability generated presence
($\beta = 0.540$) while presence's path to learning was null
($\beta = 0.069$, $p = .367$; Merchant et al.\ 560).} Parong and Mayer ran
the media comparison with content matched line for line: a narrated
slideshow beat immersive VR on the knowledge posttest in both experiments
while VR won every affective measure taken---enjoyment, motivation,
engagement, interest (``Learning Science'').\footnote{The slideshow's advantage on the knowledge
test given after the lesson was d\,=\,0.92 in the first experiment and 1.12
in the second. The same
study carries the constructive result: a generative summarizing prompt
inside VR closed the learning gap without reducing enjoyment---the deficit
is remediable by instructional design, not intrinsic to the medium.} Their
mediation follow-up closed the causal loop: the medium's effect on retention
ran through self-reported extraneous load and high-arousal positive
affect---the features that make immersive VR appealing are the ones
implicated in its learning cost (``Cognitive and Affective
Processes'' 236--37).
% [Short title normalized at assembly 2026-07-29 to match M7.1's
%  "Cognitive and Affective Processes" (MLA first-noun-phrase form);
%  per M7 header's "normalize at assembly" instruction.]
Mayer, Makransky, and Parong's review of the program's two decades states
the upshot plainly: ``[t]here is not strong evidence that students learn
academic material more deeply when it is presented in IVR than when it is
presented with conventional media'' (2236). Barrett and colleagues sharpen
the point from the psychophysics side: on one category-learning task,
learning accuracy was statistically equivalent across VR, desktop-3D, and 2D
versions; what differed was response time, fixation behavior, and---in VR
alone---attrition from physical discomfort (12--14).\footnote{The
equivalence claim is accuracy-specific, and the 2D arm is small
(n\,=\,8)---a scope limit worth carrying with the citation.}

The health-professions simulation meta-analyses locate where screen-based
practice does reliably improve learning: in skills and clinical reasoning
rather than declarative knowledge.
Kononowicz and colleagues' meta-analysis of screen-based virtual patients
found nothing added on knowledge over traditional formats (SMD\,=\,0.11,
non-significant) and a moderate-to-large advantage on skills and clinical
reasoning (SMD\,=\,0.90) (6--7). The base rate behind that split is settled:
across 609 studies of technology-enhanced simulation, pooled effects against
no intervention are large on every outcome family, and the reviewers close
the question themselves---``we question the need for further studies
comparing simulation with no intervention'' (Cook et al.\
987).\footnote{The two meta-analyses do not overlap: Cook and colleagues
pool technology-enhanced simulation broadly and explicitly exclude the
screen-based virtual patients Kononowicz and colleagues analyze. The
virtual-patient meta-analysis of that earlier era is Cook, Erwin, and
Triola (2010), a different study, not drawn on here.} Simulation teaches;
the evidence thins exactly where a world is only watched. The knowledge/
skills split names where the medium's value lives: not in what a student
sees from inside it, but in what the student is given to do there.

Read together, these two decades leave a precise explanatory hole. Presence
belongs to the person; presence does not convert into learning on its own;
learning gains concentrate where the environment gives its student something
to do. The field has measured all three results and has no construct that
joins them---no name for the difference between being placed in a world and
being taken into it as an agent whose acts complete what the world offers.
Game studies built that construct: in Gordon Calleja's \textit{In-Game: From
Immersion to Incorporation}, incorporation names the double movement in which
the virtual environment is absorbed into the player's consciousness while the
player is embodied within it as an agent. The
education literature does not import it; the bridge between the two fields
remains unbuilt in both directions. Part II rebuilt the construct on this
dissertation's own ground---rhetorical incorporation, one in substance and
dual in being, world-disclosedness and co-disclosedness---and the desktop
deployment now gives it empirical work to do. The survey data show a world
that places its students: presence attested on ordinary monitors, exactly
where the field's own findings say it can live. What the data must then
answer is whether the world completes what it offers---and that is a
question about rhetorical incorporation, not about presence, which is why the frame
Parts I--II built is the instrument this section reads with.

% --- M1.4 The asynchronous/scale rationale (batch 3, 2026-07-20)
% Quotes PDF-verified per bank §3. Additional Works Cited entries:
%   Cossio, [first names pinned at assembly], et al. "Cybersickness and
%     Discomfort from Head-Mounted Displays Delivering Fully Immersive
%     Virtual Reality: A Systematic Review." Nurse Education in Practice,
%     vol. 85, 2025, 104376.
%   Makransky, Guido, et al. "Equivalence of Using a Desktop Virtual
%     Reality Science Simulation at Home and in Class." PLOS ONE, vol. 14,
%     no. 4, 2019, e0214944.
%   Radianti, Jaziar, et al. "A Systematic Review of Immersive Virtual
%     Reality Applications for Higher Education." Computers & Education,
%     vol. 147, 2020, 103778. [full subtitle pinned at assembly]

Whether an at-home, unsupervised deployment gives up anything a supervised
classroom would provide has been tested directly. Makransky, Mayer, Veitch,
and colleagues had 112 students complete the same desktop science simulation
either at home or in class: learning outcomes did not differ (d\,=\,0.14),
nor did intrinsic motivation (d\,=\,0.02) or self-efficacy (d\,=\,0.09),
and the null pattern survived adjustment for prior knowledge and goal
orientation (``Equivalence'' 8--9).\footnote{One interaction survived
adjustment: mastery orientation predicted self-efficacy more strongly in
the classroom group (``Equivalence'' 9--10). The sample is self-selected
(112 of 289 eligible students consented), a limit the authors flag
themselves (11--12).} Their interpretation extends Clark's old verdict on
media to physical context: the instructional method---the simulation
itself---drives the learning, and felt presence is proposed, though not
tested, as what renders the location unimportant (10). The scope of the
design deserves equal attention. Both groups worked with the identical
desktop simulation---a Labster-built virtual biology laboratory---and only
the setting changed: one group at home, unsupervised, the other in a
classroom with a teacher present. No group received a different medium.
The study therefore answers one question and leaves a second untouched.
It shows that moving a desktop simulation out of the classroom and into
the student's own home, on the student's own schedule, costs nothing
measurable in learning, motivation, or self-efficacy---the fact an
asynchronous online course needs established. It does not show that an
interactive 3D simulation teaches better than a flat video of the same
content, because no video condition existed in the design. That second
question is the one our own deployment poses: every term, part of the
class entered the interactive desktop world while another part watched
the same operating-room footage on YouTube instead, and the survey data
speak directly to what separated those two experiences.

The field's own deployment history tells the complementary story. Radianti
and colleagues, reviewing 38 immersive-VR applications in higher education,
found that ``the majority of articles indicated that VR for education is
still in its experimental stages---prototyping and testing with students,''
with evaluation skewed toward usability rather than learning and the
maturity level itself ``a barrier for its adoption in regular teaching
activities'' (22).\footnote{The review's inclusion criterion is
head-mounted immersive VR only---desktop VR is excluded by design---so the
prototype-stage verdict describes precisely the hardware route, not the
desktop one (24).} A virtual world in routine use every term, across a
university system, is what that literature documents almost never
happening; the desktop build is what makes the routine part attainable.
The headset route also carries a physical cost the desktop route does not
levy. Cossio and colleagues' review of 25 HMD studies reports a mixed
rather than uniform picture---twelve studies with inconsistent or very low
cybersickness scores, five with high scores, oculomotor disturbance the
most common symptom---but where head-mounted delivery was compared directly
against desktop or non-immersive formats, the worse outcomes fell on the
headset side, ``likely due to its more immersive nature and the potential
for greater sensory conflict'' (14).\footnote{The aggregate is a reportable
moderator, not a uniform verdict: Huber and colleagues (2018) found
cybersickness \textit{less} prevalent in fully immersive VR (3.3\%) than in
augmented reality (6.6\%) (Cossio et al.\ 14). The same cost surfaces as
the VR-only discomfort attrition in Barrett and colleagues' three-arm
comparison.} Three findings from this literature bear directly on the
shape of our course. Students learn as much from a desktop simulation
completed at home as from the same simulation completed in a supervised
classroom (``Equivalence'' 8--9). Headset VR in higher education remains,
by the field's own account, a prototype technology---built and tested with
students, only rarely made part of a regular course (Radianti et al.\ 22).
Headsets carry a documented cybersickness cost that ordinary monitors do
not (Cossio et al.\ 14). Our course is fully online and asynchronous,
serving undergraduate biomedical engineering students across the
University of California's campuses, each on their own machine and their
own schedule. For a course of that shape, the desktop deployment is not a
diminished substitute for headset delivery. The home-versus-classroom
evidence says our students lose nothing by working where and when they
want. The deployment reviews say the headset alternative has almost never
been made routine at any scale. The cybersickness literature says the
monitor spares them a physical cost besides. The desktop build is what
the field's own evidence recommends for the course we actually teach.

% --- M1.6 Learning together in shared virtual space (batch 4, 2026-07-20)
% Quotes PDF-verified per bank §4. Ledger: seam DOCUMENTED here, argued in
%   M0.5 (other session); Biocca/Terry & Doolittle at statement level (argued
%   use reserved for M4.3); Kreijns Pitfall 1 quoted, Pitfall 2 glossed (full
%   deployment reserved for M5.2). Additional Works Cited entries:
%   Biocca, Frank, Judee K. Burgoon, and [initials at assembly] Harms.
%     "Toward a More Robust Theory and Measure of Social Presence: Review
%     and Suggested Criteria." Presence: Teleoperators and Virtual
%     Environments, vol. 12, no. 5, 2003, pp. 456-80.
%   De Back, Tycho, et al. "Learning in Immersed Collaborative Virtual
%     Environments: Design and Implementation." [venue + pages pinned at
%     assembly], 2021.
%   Garrison, D. Randy, Terry Anderson, and Walter Archer. "Critical
%     Inquiry in a Text-Based Environment: Computer Conferencing in Higher
%     Education." The Internet and Higher Education, vol. 2, no. 2-3,
%     1999, pp. 87-105.
%   Gunawardena, Charlotte N., and Frank J. Zittle. "Social Presence as a
%     Predictor of Satisfaction within a Computer-Mediated Conferencing
%     Environment." American Journal of Distance Education, vol. 11,
%     no. 3, 1997, pp. 8-26.
%   Kreijns, Karel, Paul A. Kirschner, and Wim Jochems. "Identifying the
%     Pitfalls for Social Interaction in Computer-Supported Collaborative
%     Learning Environments: A Review of the Research." Computers in Human
%     Behavior, vol. 19, 2003, pp. 335-53.
%   Richardson, Jennifer C., and Karen Swan. "Examining Social Presence in
%     Online Courses in Relation to Students' Perceived Learning and
%     Satisfaction." [venue pinned at assembly], 2003.
%   Richardson, Jennifer C., et al. "Social Presence in Relation to
%     Students' Satisfaction and Learning in the Online Environment: A
%     Meta-Analysis." [venue pinned at assembly], 2017.
%   Short, John, Ederyn Williams, and Bruce Christie. The Social
%     Psychology of Telecommunications. Wiley, 1976.
%   Terry, Krista, and Peter Doolittle. "Re-examining Social Presence:
%     Implications for Digital Pedagogies." Technology, Instruction,
%     Cognition and Learning, vol. 11, 2019, pp. 121-39.
%   Tu, Chih-Hsiung. [title/venue pinned at assembly], 2002.
%   van der Meer, [first names + byline form pinned at assembly], et al.
%     [title pinned], Frontiers in Virtual Reality, vol. 4, 2023, 1159905.

The section's second literature studies a different presence: not the sense
of being in a place but the sense of being with someone. Social presence
entered the research vocabulary in 1976, in Short, Williams, and Christie's
social psychology of telecommunications, naming how fully a communication
medium lets its users register one another as persons---decades before
rendered worlds existed to be present in. Biocca, Burgoon, and Harms,
surveying the construct's first quarter-century, define its core as ``the
sense of being with another'' (456) and set out the criteria a robust
theory of it must meet. Their sharpest verdict falls on the field's own
instrument. The inherited scale was built to rate media channels, not to
describe interactions, and the authors press the question that follows:
``is social presence just an attitude towards a medium?''
(469).\footnote{The scale's provenance explains its orientation: it
descends from studies funded by the UK Post Office, the Department of
Transportation, and General Electric---organizations commissioning
ratings of their communication channels (Biocca et al.\ 469).} A
construct meant to describe an interaction had been operationalized as an
opinion about equipment.

What the construct predicts, in its home tradition of online education, is
substantial. In Gunawardena and Zittle's origin study of a text-only
academic conference, social presence alone accounted for roughly 60
percentage points of the variance in learner satisfaction
(18--19).\footnote{Their closing design recommendation cites Johansen,
Vallee, and Spangler's suggestion that social presence can ``be
cultured''---built up by design and practice rather than fixed by the
medium (Gunawardena and Zittle 23).} Richardson and Swan, adapting the
same scale, found social presence correlating .68 with perceived learning
(73--74)---a second result from the same research program, extending
rather than independently confirming the first.\footnote{The correlation
is strongest in class discussions (.83) and group projects (.80), weakest
in lectures, notes, and readings (.40) (Richardson and Swan
74--75).} The meta-analytic aggregate holds the pattern at scale: across
26 effect sizes, social presence correlates 0.56 with satisfaction and
0.51 with perceived learning (Richardson et al.\
410--11).\footnote{Both correlations are moderated by course length,
discipline, and by which scale was used---no two social-presence scales
are equivalent (Richardson et al.\ 410--12, 415). The outcome throughout
the meta-analysis is self-reported perceived learning, never graded
performance (411); Tu's decomposition of the construct returned
reliabilities he himself calls ``medium-low'' and an instrument offered
as ``a first step'' (Tu 43--44).} Terry and Doolittle press hardest on
what the construct has become: the Community-of-Inquiry tradition
entangled social presence with the teaching and cognitive presences
around it, and interactivity keeps being mistaken for it. The repair they
propose is to identify it ``as a psychological state, as opposed to one
that is mediated through, or as a by-product of the use of a specific
technology'' (121).\footnote{Five factors contribute to it on their
account: group cohesion, co-presence and the salience of others,
affective association, immediacy, and intimacy (Terry and Doolittle
130--31).} That relocation is this section's load-bearing result: social
presence, like spatial presence in the first literature, belongs to persons.
It also sets the construct on ground Heidegger prepared: being-with-others is
not an addition to a self-standing subject but a constitutive structure of
human existence---the world is always already one shared with others---and
the distinctive enactment of that structure is
speaking-with-one-another.\footnote{``The world of Dasein is a
\textit{with-world} [\textit{Mitwelt}]. Being-in is \textit{Being-with}
Others. Their Being-in-themselves within-the-world is \textit{Dasein-with}
[\textit{Mitdasein}]'' (BT 155; italics Heidegger's). The Marburg lectures on
Aristotle name the enactment: ``The distinctive mode of being in
being-with-one-another is in speaking-with-one-another. To set forth the
possibilities of being-with-one-another is the \gk{ἔργον} of rhetoric''
(\textit{Basic Concepts} 91--92).} A design that opens room for
speaking-with-others therefore does more than decorate a channel: it
furnishes the enactment in which others are there at all. It is something a
design can activate, not something a medium contains.

The same tradition states what activating it requires. Garrison, Anderson,
and Archer locate social presence in the communication context rather than
the medium and name teaching presence ``[t]he binding element in creating
a community of inquiry'' (96)---when text-based education fails, their
diagnosis is a deficit of teaching presence, not of bandwidth. Kreijns,
Kirschner, and Jochems give the tradition its bluntest warning, naming the
first of two pitfalls of computer-supported collaborative learning:
``taking for granted that social interaction automatically takes place
just because an environment makes it technologically possible'' (336)---
the second being the restriction of what interaction does occur to
task-focused exchanges. Providing the channel is not producing the
contact; sociability must be designed for. The findings extend into
rendered worlds. De Back, Tinga, and Louwerse's immersed collaborative
environments produced learning gains with an inverse group-size
gradient---smaller groups, larger gains (5375--76).\footnote{Cited for
collaborative-learning design and outcomes only; social presence does not
appear in the study's text.} Van der Meer and colleagues' review of 139
virtual-reality collaborative-learning systems found 78.4\% of them
monitor-based against 7.2\% head-mounted (6).\footnote{Among the nine
design features the review requires of such environments, verbal audio
communication is judged essential and support for socialization is named
a primary function (van der Meer et al.\ 9--11).} Collaborative
learning in virtual worlds, as actually practiced and studied, already
runs on ordinary screens.

These two literatures share their origin and almost nothing after it.
Short, Williams, and Christie stand at the head of both: the presence
tradition the VR-pedagogy field inherited and the social-presence
tradition of online education descend from the same 1976 book. Downstream,
the citation traffic stops.\footnote{The meta-analysis that pools the
social-presence evidence draws its studies from online courses, not
rendered worlds, and pools neither Biocca nor Kreijns (Richardson et al.\
415--16); the desktop and headset studies of the first literature measure
spatial presence, cognitive load, and learning outcomes, and the
social-presence instruments appear in none of them.} This section reads a
desktop world whose students meet inside it by voice, and so it must join
two bodies of evidence that have not met each other---the second
literature's vocabulary applied to an object the first literature built.
What this movement hands forward is the
base for that work: a construct that names being-with as its own register
of presence; evidence that it predicts satisfaction and perceived
learning; the warning that it must be designed for and convened, never
assumed; and the finding that it belongs to persons---available, in
principle, to any world that gives its students a way to be there
together.

% --- M1 LOCKED 2026-07-20 (user). All four movements approved: M1.1 rev 3
%   -> merged M1.2/3 -> M1.4 rev -> M1.6 rev 2 (footnote rebalance
%   approved). DO NOT EDIT until the newest batch of sources has been
%   ingested (archon + snapshot rebuild), indexed (corpus/index entry
%   additions), and analyzed for useful/necessary material — then a
%   supplement pass on M1 under the same bank/verification workflow.
%   Deferred to unlock/assembly: venue/name pins in header comments;
%   Works Cited consolidation; ledger final check before M4.3/M5.2
%   drafting in the main session.
